\chapter{สรุปและข้อเสนอแนะ}
\label{Ch:Conclusion}

บทนี้นำเสนอผลสรุปของโครงงาน \textit{PyGreenSense} ซึ่งเป็นระบบตรวจจับ Green Code Smell สำหรับภาษา Python ร่วมกับส่วนขยาย Visual Studio Code โดยสรุปผลจากการพัฒนาและการประเมินระบบในบทก่อนหน้า รวมถึงอธิบายข้อจำกัดที่พบระหว่างการดำเนินงาน และเสนอแนวทางสำหรับการพัฒนาต่อยอดในอนาคต

% -----------------------------------------------------

\section{สรุปผลการดำเนินโครงงาน}

โครงงานนี้มีวัตถุประสงค์เพื่อพัฒนาเครื่องมือที่ช่วยให้ผู้พัฒนาโปรแกรมสามารถมองเห็นปัญหาเชิงโครงสร้างของโค้ดที่อาจส่งผลต่อประสิทธิภาพและความยั่งยืนของซอฟต์แวร์ โดยเน้นการตรวจจับ \textit{Green Code Smell} ในภาษา Python และนำเสนอผลลัพธ์ผ่านส่วนขยาย Visual Studio Code เพื่อให้สามารถใช้งานได้ในสภาพแวดล้อมการพัฒนาจริง

ผลลัพธ์หลักของโครงงานคือการพัฒนาไลบรารี \textit{PyGreenSense} ที่สามารถวิเคราะห์ซอร์สโค้ดภาษา Python ด้วยแนวทาง \textit{AST-based static analysis} และตรวจจับ Code Smell ที่เกี่ยวข้องกับคุณภาพของโค้ด ได้แก่ Dead Code, Long Method, God Class, Duplicated Code และ Mutable Default Arguments นอกจากนี้ ระบบยังมีการพัฒนาโมดูลสำหรับจัดเก็บประวัติการวิเคราะห์ในรูปแบบไฟล์ JSON เพื่อรองรับการเปรียบเทียบผลก่อนและหลังการปรับปรุงโค้ด

ในเชิงแนวคิด โครงงานนี้ช่วยให้เห็นภาพของการพัฒนาโค้ดที่มีแนวโน้มเป็นมิตรต่อสิ่งแวดล้อมมากขึ้น หรือ \textit{greener code} โดยไม่ได้พิจารณาเฉพาะความถูกต้องของโปรแกรมเท่านั้น แต่ยังพยายามเชื่อมโยงคุณภาพของโค้ดเข้ากับประเด็นด้านพลังงาน คาร์บอน และความยั่งยืนของซอฟต์แวร์ การพัฒนาดังกล่าวช่วยสะท้อนให้เห็นว่า Code Smell บางประเภทอาจไม่ได้เป็นเพียงปัญหาด้านการอ่านหรือการบำรุงรักษาโค้ด แต่ยังอาจเกี่ยวข้องกับต้นทุนในการประมวลผลและการใช้ทรัพยากรของซอฟต์แวร์ด้วย

อีกประเด็นสำคัญที่โครงงานนี้แสดงให้เห็นคือ เครื่องมือวิเคราะห์โค้ดในปัจจุบันจำนวนมากยังมีข้อจำกัดในการเชื่อมโยงผลการตรวจจับ Code Smell กับผลกระทบด้านพลังงานหรือคาร์บอนโดยตรง เครื่องมือทั่วไปมักเน้นการตรวจจับปัญหาด้านคุณภาพโค้ด เช่น dead code, duplicated code หรือ complexity เป็นหลัก แต่ยังไม่ได้อธิบายว่าปัญหาเหล่านี้ส่งผลต่อมิติด้าน Green Software Engineering อย่างไร ดังนั้น \textit{PyGreenSense} จึงเป็นแนวทางหนึ่งในการทดลองเชื่อมช่องว่างระหว่างการวิเคราะห์คุณภาพโค้ดและการประเมินผลกระทบด้านความยั่งยืนของซอฟต์แวร์

% -----------------------------------------------------

\section{สรุปผลการพัฒนาระบบ}

จากการพัฒนา ระบบ \textit{PyGreenSense} ประกอบด้วยองค์ประกอบหลัก ได้แก่ AST Parser, ชุดกฎสำหรับตรวจจับ Green Code Smell, Metric and Analysis Module, VS Code Extension Module และ Reporting Module โดยแต่ละส่วนทำงานร่วมกันเพื่อให้สามารถวิเคราะห์โค้ด ตรวจจับปัญหา จัดเก็บผลลัพธ์ และแสดงผลให้ผู้ใช้เข้าใจได้ง่ายขึ้น

ในส่วนของไลบรารีตรวจจับ ระบบสามารถใช้ AST ของภาษา Python เพื่อวิเคราะห์โครงสร้างของโปรแกรมในระดับฟังก์ชัน คลาส ตัวแปร และคำสั่งต่าง ๆ ได้อย่างเป็นระบบ กฎการตรวจจับที่พัฒนาขึ้นสามารถระบุปัญหาที่พบได้ เช่น ฟังก์ชันที่ยาวเกินไป คลาสที่มีความรับผิดชอบมากเกินไป โค้ดที่ซ้ำกัน โค้ดที่ไม่ได้ใช้งาน และการใช้ค่าเริ่มต้นของพารามิเตอร์ที่เป็น mutable object

ในส่วนของมาตรวัด ระบบได้เพิ่มแนวคิดการคำนวณ \textit{Cosmic Function Points (CFP)} เพื่อประมาณขนาดเชิงฟังก์ชันของซอฟต์แวร์จากซอร์สโค้ด และนำแนวคิด \textit{Software Carbon Intensity (SCI)} มาประยุกต์ใช้ในรูปแบบ SCI per CFP เพื่อพิจารณาปริมาณคาร์บอนต่อหน่วยงานที่ซอฟต์แวร์ทำจริง นอกจากนี้ ระบบยังมีตรรกะสำหรับเปรียบเทียบผลก่อนและหลังการปรับปรุงโค้ด เช่น Impact Logic และ Rule-based Estimated Saving เพื่อใช้วิเคราะห์แนวโน้มผลกระทบของการลด Code Smell

ในส่วนของ VS Code Extension ระบบถูกออกแบบให้ช่วยให้ผู้ใช้สามารถเรียกใช้งานการวิเคราะห์ได้สะดวกขึ้นจากภายในเครื่องมือพัฒนาโค้ด โดยมีแนวคิดในการจัดการ environment ของไลบรารี การเรียกใช้งานระบบตรวจจับ และการแสดงผลลัพธ์ในรูปแบบที่ผู้ใช้สามารถนำไปใช้ประกอบการตัดสินใจปรับปรุงโค้ดได้

% -----------------------------------------------------

\section{สรุปผลการประเมิน}

ผลการประเมินแสดงให้เห็นว่า \textit{PyGreenSense} สามารถตรวจจับ Green Code Smell ตามกฎที่กำหนดไว้ได้ และสามารถนำข้อมูลผลลัพธ์ไปใช้ประกอบการวิเคราะห์ก่อนและหลังการปรับปรุงโค้ดได้ในระดับหนึ่ง โดยเฉพาะการเปรียบเทียบจำนวน Code Smell ที่ลดลงหลังการ refactor และการจัดเก็บข้อมูลใน History Log เพื่อใช้ติดตามผลลัพธ์ในระยะยาว

อย่างไรก็ตาม ผลการประเมินยังสะท้อนข้อจำกัดสำคัญของการทดลองในปัจจุบัน กล่าวคือ ชุดโค้ดที่ใช้ในการทดสอบยังมีจำนวนและความหลากหลายจำกัด ทำให้ผลลัพธ์ด้านพลังงานและคาร์บอนยังไม่สามารถสรุปได้อย่างครอบคลุมในทุกสถานการณ์ โดยเฉพาะในกรณีที่ source code มีขนาดเล็กหรือมี workload ไม่มากพอ ค่าพลังงานที่วัดได้อาจมีความผันผวนจากสภาพแวดล้อมของเครื่องที่ใช้ทดสอบ และอาจทำให้การตีความผลกระทบของการ refactor ต่อการลดคาร์บอนยังไม่ชัดเจน

นอกจากนี้ ผลกระทบของ Code Smell แต่ละประเภทต่อการใช้พลังงานยังมีข้อมูลเชิงประจักษ์ค่อนข้างน้อย แม้ระบบจะสามารถตรวจจับและลดจำนวน Code Smell ได้ แต่ยังไม่สามารถสรุปได้แน่ชัดว่า Code Smell แต่ละประเภทส่งผลต่อพลังงานหรือคาร์บอนมากน้อยเพียงใดในทุกกรณี ตัวอย่างเช่น การลด Dead Code หรือ Long Method อาจช่วยให้โค้ดอ่านง่ายและบำรุงรักษาได้ดีขึ้นอย่างชัดเจน แต่ผลด้านพลังงานอาจขึ้นอยู่กับลักษณะการทำงานจริงของโปรแกรม ขนาดข้อมูลที่ใช้ทดสอบ และจำนวนครั้งที่โค้ดส่วนนั้นถูกเรียกใช้งาน

ดังนั้น ข้อสรุปสำคัญจากการประเมินคือ การลด Code Smell สามารถช่วยปรับปรุงคุณภาพของโค้ดและทำให้แนวทางการพัฒนาโค้ดมีความเป็นระบบมากขึ้น แต่ผลกระทบด้านพลังงานและคาร์บอนยังควรถูกตีความในฐานะผลลัพธ์เชิงประมาณการหรือแนวโน้มเบื้องต้น มากกว่าการยืนยันผลเชิงสรุปในทุกสถานการณ์

% -----------------------------------------------------

\section{ข้อจำกัดของโครงงาน}

แม้ว่าโครงงานนี้จะสามารถพัฒนาเครื่องมือตรวจจับ Green Code Smell และเชื่อมโยงกับมาตรวัดด้านคาร์บอนได้ในระดับหนึ่ง แต่ยังมีข้อจำกัดที่ควรพิจารณา ดังนี้

\begin{itemize}
    \item {\boldfont ข้อจำกัดด้านชุดโค้ดทดสอบและ Bias ของข้อมูล}
    ชุดโค้ดที่ใช้ยังมีขนาดและความหลากหลายจำกัด และบางส่วนเป็น {\italicfont mock code} ที่มีลักษณะเรียบง่าย เช่น {\italicfont dead code} หรือ {\italicfont mutable default argument} แบบตรงไปตรงมา รวมถึงการ refactor ที่ไม่สม่ำเสมอ ทำให้ไม่สามารถสะท้อนสถานการณ์จริงได้ครบถ้วน และอาจเกิด {\italicfont bias} ที่ส่งผลให้ผลลัพธ์ยังไม่สามารถ {\italicfont generalize} ไปยังระบบจริงได้

    \item {\boldfont ข้อจำกัดด้านผลกระทบของ Code Smell แต่ละประเภท}
    ข้อมูลเชิงประจักษ์เกี่ยวกับผลกระทบของ Code Smell แต่ละประเภทต่อการใช้พลังงานยังมีไม่มากพอ ทำให้การวิเคราะห์ผลกระทบของแต่ละ smell ยังเป็นการประมาณจากข้อมูลที่ระบบตรวจพบ เช่น จำนวนบรรทัดที่ถูกลดลง จำนวน smell ที่ลดลง หรือค่าคาร์บอนที่เปลี่ยนแปลงหลังการทดสอบ

    \item {\boldfont ข้อจำกัดด้านการวัดพลังงานและคาร์บอน}
    ผลการวัดพลังงานและคาร์บอนอาจได้รับผลกระทบจากสภาพแวดล้อมของเครื่องที่ใช้ทดสอบ เช่น background process, hardware, operating system และลักษณะการรันโปรแกรมในแต่ละครั้ง ทำให้ค่าที่ได้อาจมีความผันผวน

    \item {\boldfont ข้อจำกัดด้านขอบเขตของ Static Analysis}
    ระบบใช้การวิเคราะห์จากโครงสร้าง AST เป็นหลัก จึงสามารถตรวจจับปัญหาเชิงโครงสร้างของโค้ดได้ดีในระดับหนึ่ง แต่ยังมีข้อจำกัดในการวิเคราะห์พฤติกรรมที่เกิดขึ้นจริงระหว่าง runtime เช่น input ที่เปลี่ยนแปลงตามสถานการณ์ หรือโค้ดที่ถูกเรียกใช้งานผ่าน dynamic behavior

    \item {\boldfont ข้อจำกัดด้านการใช้งานของผู้เริ่มต้น}
    แม้ระบบจะช่วยแสดงผล Code Smell ได้ แต่การอธิบายผลลัพธ์และการแนะนำวิธีแก้อาจยังเข้าใจยาก และอาศัยความเข้าใจค่อนข้างสูงสำหรับผู้เริ่มต้นเขียนโค้ด

    \item {\boldfont ข้อจำกัดของ SCI (Software Carbon Intensity) ในบางกรณี}
    แม้จะลด LOC, CFP หรือจำนวน {\italicfont code smell} ลงอย่างมาก แต่ค่าคาร์บอนหรือ SCI อาจไม่ได้ลดลงตาม หรืออาจเพิ่มขึ้น เนื่องจาก SCI ขึ้นอยู่กับพฤติกรรม {\italicfont runtime} และปัจจัยแวดล้อมอื่น ๆ ทำให้ไม่สามารถใช้เป็นตัวชี้วัดผลของการ refactor ได้โดยตรงในทุกกรณี
\end{itemize}

% -----------------------------------------------------

\section{ข้อเสนอแนะสำหรับการพัฒนาต่อ}

จากผลการดำเนินงานและข้อจำกัดที่พบ สามารถเสนอแนวทางพัฒนาต่อในอนาคตได้ดังนี้

\begin{itemize}
    \item {\boldfont การปรับปรุงมาตรฐานและความน่าเชื่อถือของการวัด (Standardization \& Reliability)}
    ควรพัฒนา {\italicfont baseline}, {\italicfont threshold} และ {\italicfont Data Catalog} สำหรับผลกระทบของ Code Smell ต่อพลังงานและคาร์บอน เพื่อให้สามารถเปรียบเทียบผลได้อย่างเป็นมาตรฐานและลดความคลาดเคลื่อนของการวิเคราะห์

    \item {\boldfont การพัฒนา Refactoring Suggestion และ Editor Integration}
    ควรพัฒนาระบบแนะนำแนวทางแก้ไข Code Smell และเชื่อมต่อกับ {\italicfont code editor} เช่น VS Code Extension เพื่อช่วยให้ผู้ใช้สามารถปรับปรุงโค้ดได้แบบ {\italicfont real-time} และใช้งานได้สะดวกยิ่งขึ้น

    \item {\boldfont การขยายกฎการตรวจจับ Code Smell}
    ควรเพิ่มจำนวนและความหลากหลายของ {\italicfont code smell rules} รวมถึงรองรับ {\italicfont pattern} ที่พบในโปรเจกต์จริงมากขึ้น เพื่อให้ครอบคลุมปัญหาด้านคุณภาพโค้ดได้ดียิ่งขึ้น

    \item {\boldfont การพัฒนาการแสดงผลและ Visualization (Gamification)}
    ควรออกแบบการแสดงผลให้น่าสนใจมากขึ้น เช่น การใช้ {\italicfont score}, {\italicfont badge} หรือระดับ ({\italicfont level}) เพื่อช่วยให้ผู้ใช้เข้าใจผลลัพธ์ได้ง่ายและสร้างแรงจูงใจในการเขียนโค้ดที่มีประสิทธิภาพมากขึ้น

    \item {\boldfont การเพิ่มประสิทธิภาพสำหรับโปรเจกต์ขนาดใหญ่ (Scalability)}
    ควรปรับปรุงประสิทธิภาพของระบบให้สามารถวิเคราะห์ {\italicfont source code} ขนาดใหญ่ได้รวดเร็วและมีประสิทธิภาพมากขึ้น รวมถึงรองรับโครงสร้างโปรเจกต์ที่ซับซ้อน

    \item {\boldfont การเพิ่มชุดทดสอบและข้อมูลจากการใช้งานจริง}
    ควรรวบรวม {\italicfont source code} และ {\italicfont log} จากการใช้งานจริงเพื่อนำมาวิเคราะห์และลด {\italicfont bias} ของ {\italicfont mock data} รวมถึงช่วยให้ผลลัพธ์มีความสมจริงและน่าเชื่อถือมากขึ้น
\end{itemize}

% -----------------------------------------------------

\section{สรุปภาพรวม}

โดยสรุป โครงงาน \textit{PyGreenSense} สามารถพัฒนาเครื่องมือตรวจจับ Green Code Smell สำหรับภาษา Python และแสดงให้เห็นแนวทางเบื้องต้นในการเชื่อมโยงคุณภาพของโค้ดเข้ากับประเด็นด้านพลังงานและคาร์บอนของซอฟต์แวร์ได้สำเร็จ โครงงานนี้ช่วยให้เห็นภาพว่าการพัฒนาโค้ดที่ดีไม่ได้จำกัดอยู่เพียงความถูกต้องหรือความง่ายต่อการบำรุงรักษาเท่านั้น แต่ยังสามารถขยายไปสู่แนวคิดการพัฒนาโค้ดที่คำนึงถึงความยั่งยืนได้ด้วย

อย่างไรก็ตาม ผลลัพธ์ด้านพลังงานและคาร์บอนยังมีข้อจำกัดจากชุดทดสอบและข้อมูลผลกระทบของแต่ละ Code Smell ที่ยังมีจำนวนน้อย ดังนั้นการพัฒนาต่อควรมุ่งเน้นการเก็บข้อมูลเชิงประจักษ์เพิ่มเติม การสร้าง Data Catalog ของผลกระทบ Code Smell ในภาษา Python และการนำข้อมูลจากการใช้งานจริงของ Extension และ Library มาวิเคราะห์ เพื่อให้สามารถพัฒนาเครื่องมือที่ช่วยผู้ใช้เขียนโค้ดได้ดีขึ้น มีประสิทธิภาพมากขึ้น และตระหนักถึงผลกระทบด้านพลังงานของซอฟต์แวร์มากขึ้นในอนาคต
